\scanpage{22}

\uline{\unsure{Remarques|Remarque}}\quad 1) \={M} si $\psi$ est \deleted{\unsure{gén.} étale}\inserted{\uline{birationnel}}, \uline{et} $\mathfrak h$ provient
d'un sous-groupe $H$ de $G$ lisse sur $k$, \deleted{\unsure{la même|aussi}}
la condition (i) n'est pas néc. vérifiée.

\uline{Exemple}\quad $G=\mathrm{Sl}(2,k)$\quad $k$ car. 2, Alors
$\mathfrak g$ est \uline{nilpotent} donc \deleted{\unsure{les}}la \inserted{seule} sous-algèbre de
Cartan \deleted{\unsure{de sont|de son}} est $\mathfrak g$ elle \={m}. Soit
$B$ Borel, alors \deleted{$\mathscr L$} $\mathfrak h=\operatorname{Lie}(B)$, on sait
que $\bigcup_{g\in G(k)}\operatorname{ad}(g)\mathscr L=\mathfrak g$ \deleted{\unsure{, et}} \unsure{(cf)|(sic)}, \deleted{\unsure{et|donc}}
\deleted{\unsure{est}} donc $\psi$ dominant. D'ailleurs
$M=\mathfrak g\quad B\subset N\ne G$, donc $\dim N=\dim B=\operatorname{rang}\mathfrak h$,
donc $\dim X=\dim G$, \deleted{D'ailleurs on a} \inserted{donc $\psi$ gén. quasi fini.}

\deleted{aussi pour $\underline{\forall}\ a\in\mathfrak h$, $M=M_a=\mathfrak g$\quad (car $M\subset M_a$)}

D'ailleurs, une analyse plus poussée montre
\deleted{donc $\varphi$ \uline{lisse} en tout point \unsure{gén.} de \unsure{$X$}}
\deleted{$\psi$ étale en tout point}

\diagram{\begin{tikzcd}[ampersand replacement=\&,row sep=small,column sep=small]
\& \deleted{N=M_{a}} \& \\
\deleted{\unsure{G/N|G_N}} \& \deleted{\unsure{M=M_a|M=\mathfrak g}} \& \deleted{\begin{pmatrix}a&b\\0&-a\end{pmatrix}}\\
\deleted{\unsure{G|\mathscr L}\quad B\subset G} \& \&
\end{tikzcd}}

ici que $\psi$ \uline{est}
\deleted{\unsure{quasi-ét.}} \uline{birationnel}
\deleted{\unsure{un système de B et B$^-$|$B\cap B$}}

\reversemarginpar
\margin{\setlength{\fboxsep}{1pt}\fbox{\parbox{0.94\marginparwidth}{\raggedright\scriptsize\uline{N.B.}\quad $\varphi$ lisse en $(\dot e,a)$\\
$\Updownarrow$\\
$\operatorname{ad}(a)_{\mathfrak g/\mathfrak h}$ injectif\\
i.e. $[x,a]\in\mathfrak h\Rightarrow x\in\mathfrak h$\\
$M_a=\mathfrak h$\\
$\Downarrow$\\
$\mathfrak h=M$}}}
\normalmarginpar

2) \={M} si $\psi$ \deleted{birationnel}\inserted{dominant} et $\mathfrak h=M$
la condition (\unsure{x|i|xi}) n'est pas
néc. vérifiée. Prenons
$G=\mathrm{Pl}(2,k)$, $k$ car 2, $\mathfrak g$ a une
base $X,Y,H$ avec
$[X,Y]=2H=0,\quad[H,X]=X,\quad[H,Y]=Y$.

% archive sheet number: 21
